% Part: first-order-logic
% Chapter: axiomatic-deduction
% Section: deduction-theorem-quantifiers

\documentclass[../../../include/open-logic-section]{subfiles}

\begin{document}

\olfileid{fol}{axd}{ddq}

\olsection{పరిమాణీకరణికాలతో నిగమన సిద్ధాంతం}

\begin{thm}[నిగమన సిద్ధాంతం]
\ollabel{thm:deduction-thm-q} $\Gamma \cup \{!A\} \Proves !B$
అయితే, $\Gamma \Proves !A \lif !B$.
\end{thm}

\begin{proof}
$\Gamma \cup \{!A\}$ నుంచి $!B$కు గల \tetoken{వ్యుత్పత్తి}{!!{derivation}} పొడవుపై మళ్లీ
ఆగమనం చేస్తాం.

ఆగమన ప్రాతిపదిక నిరూపణ,~\olref[ded]{thm:deduction-thm} నిరూపణలో
ఉన్నదానితో సరిసమానం.

ఆగమన దశలో, $\Gamma \cup \{!A\}$ నుంచి $!B$కు గల \tetoken{వ్యుత్పత్తి}{!!{derivation}},
ఒక నిగమన నియమం ద్వారా సమర్థించబడిన~$!B$ దశతో ముగుస్తుందని మళ్లీ
అనుకుందాం. ఆ నియమం మోడస్ పోనెన్స్ అయితే,
\olref[ded]{thm:deduction-thm} నిరూపణలోలాగే సాగుతాం. ఆ నిగమన నియమం
\QR{} అయితే, $!B \ident !C \lif \lforall[x][!D(x)]$ అని, అలాగే
$!C \lif !D(a)$ రూపంలోని ఒక \tetoken{సూత్రం}{!!a{formula}} ఆ \tetoken{వ్యుత్పత్తిలో}{!!{derivation}} ముందుగా
కనిపిస్తుందని మనకు తెలుసు; ఇక్కడ $a$ అనేది~$!C$, $!A$, $\Gamma$లలో
కనిపించదు. అందువల్ల
\begin{align*}
  \Gamma \cup \{!A\} & \Proves !C \lif !D(a),\\
  \intertext{అలాగే ఆగమన పరికల్పన వర్తిస్తుంది; అంటే,}
    \Gamma & \Proves !A \lif (!C \lif !D(a)).\\
  \intertext{ఇంకా}
  & \Proves (!A \lif (!C \lif !D(a))) \lif ((!A \land !C) \lif !D(a))\\
  \intertext{మరియు మోడస్ పోనెన్స్ వల్ల}
  \Gamma & \Proves (!A \land !C) \lif !D(a).\\
  \intertext{ఐగెన్ చరరాశి షరతు ఇంకా నెరవేరుతుంది కనుక, ఈ \tetoken{వ్యుత్పత్తికి}{!!{derivation}} \QRతో సమర్థించిన ఒక దశను చేర్చి}
    \Gamma & \Proves (!A \land !C) \lif \lforall[x][!D(x)].\\
    \intertext{ఇంకా మనకు}
    & \Proves ((!A \land !C) \lif \lforall[x][!D(x)]) \lif (!A \lif (!C \lif \lforall[x][!D(x)])),\\
    \intertext{కాబట్టి మోడస్ పోనెన్స్ వల్ల}
    \Gamma & \Proves !A \lif (!C \lif \lforall[x][!D(x)]),
\end{align*}
అంటే, $\Gamma \Proves !A \lif !B$.
\sourcecorrection{OLTEAXD-004}{ఆంగ్ల మూలం రెండవ మెటా-నిబంధనలోని
పర్యవసానాన్ని మూసే కుడి కుండలీకరణ చిహ్నాన్ని వదిలింది. మొత్తం
$(!A \lif (!C \lif \lforall[x][!D(x)]))$ పర్యవసానం మూసేలా
సరిచేశాం.}
\sourcecorrection{OLTEAXD-005}{ఈ సందర్భంలో $!B \ident !C \lif
\lforall[x][!D(x)]$; నిరూపించాల్సిన నిగమన-సిద్ధాంత ముగింపు
$\Gamma \Proves !A \lif !B$. ఆంగ్ల మూలం చివర పొరపాటుగా
$\Gamma \Proves !B$ అని రాసింది; లక్ష్యానికి సరిపోయేలా సరిచేశాం.}

$!B$ను \QR{} నియమం సమర్థించినా, అది $\lexists[x][!D(x)] \lif !C$
రూపంలో ఉండే సందర్భాన్ని సాధనగా వదిలేస్తాం.
\end{proof}

\begin{prob}
  \olref[fol][axd][ddq]{thm:deduction-thm-q} నిరూపణను పూర్తి చేయండి.
\end{prob}

\end{document}
